\chapter{Extensiones nulas y clausura material}
\label{chap:realizaciones-nula-material}
\index{fotón!interfaz física}\index{electrón!clausura material}
\index{extensión nula}\index{cono causal partido}

\section{\(4\) realizaciones del estado material común: nula, masiva, compuesta y topológica}
\label{sec:cuatro-modos-realizacion}

Una cifra aislada determina una coordenada, pero no identifica la estructura
que la produce, la memoria que conserva ni la operación que la convierte en
magnitud. HMT se sitúa antes de esa proyección: su objeto primario es una
extensión compatible que conserva las leyes aritméticas, la orientación
TRIT, los cocientes, la fase y la memoria de frontera. Mecánica Dimensional
especifica después qué componente se realiza como magnitud, cuál permanece
como memoria y qué transformación produce el observable.

Los desarrollos de esta parte distinguen cuatro modos por su operación de
cierre:
\begin{enumerate}
 \item una extensión abierta y nula transporta una perturbación entre
 fronteras;
 \item una extensión cerrada y persistente adquiere norma interior y define
 el candidato material;
 \item una familia de extensiones ponderada por acción define un ensamble
 térmico;
 \item una conexión que compara secciones en puntos distintos proporciona el
 codominio de la realización gravitatoria.
\end{enumerate}
La fuente común no identifica los cuatro codominios. Los mapas físicos
correspondientes se mantienen discontinuos:
\[
 \mathfrak R_\gamma:
 \Omega_{\HMT}^{\rm surv}\dashrightarrow\mathcal H_{\rm rad},
 \qquad
 \mathfrak R_m:
 \Omega_{\HMT}^{\rm surv}\dashrightarrow\mathcal H_{\rm mat},
\]
\[
 \mathfrak R_\Theta:
 \mathcal P(\Omega_{\HMT}^{\rm surv})\dashrightarrow\mathcal S_{\rm th},
 \qquad
 \mathfrak R_{\rm grav}:
 \Omega_{\HMT}^{\rm surv}\dashrightarrow\mathsf{Cartan}_{1,3}.
\]
Su unidad reside en la genealogía; su diferencia, en la lectura y en las
estructuras adicionales exigidas por cada realización.

\section{Propagación nula entre fronteras}
\label{sec:propagacion-nula-fotonica}

Sea \(B\) una base de transporte. Una \emph{sección nula} de signo
\(\pm\) es una aplicación
\[
 s_\pm:B\longrightarrow\mathcal A_{\rm sp}P_\pm
\]
tal que \(s_\pm(x)\in\mathcal A_{\rm sp}P_\pm\) para todo \(x\in B\).
Por la \cref{prop:descomposicion-espectral-partida},
\[
 N_{\rm sp}(s_\pm(x))=0
 \qquad(x\in B).
\]
Si \(\gamma\) es una ruta orientada entre dos fronteras,
\[
 \partial\gamma=x_{\rm abs}-x_{\rm em},
\]
la interfaz fotónica propuesta asigna a \(\gamma\) una sección nula
extendida. El adjetivo \emph{extendida} indica que su identidad pertenece al
transporte completo y no a un punto aislado.

La ruta conjugada se obtiene por inversión,
\[
 \gamma^\vee=C_{\rm rev}\gamma,
 \qquad
 \partial\gamma^\vee=x_{\rm em}-x_{\rm abs}.
\]
Relativamente al lector de orientación declarado, el par
\((\gamma,\gamma^\vee)\) representa las lecturas avanzada y retardada. Las
dos piernas se ordenan secuencialmente o se alojan en fibras distintas. No se
suman automáticamente dentro de una misma fibra: si se activan
simultáneamente las dos componentes conjugadas, la identidad exacta
\[
 \boxed{N_{\rm sp}(uP_++vP_-)=uv}
\]
muestra que, para \(u,v>0\), la sección abandona la frontera nula y entra en
el interior positivo.

\subsection{Longitud, ida y retorno}

Fijada la carta dimensional de velocidad,
\(\widehat c=e^{-\mathcal E}\), y una base positiva \(u_L\) de la recta de
longitud, se define
\[
 L_{\rm int}(\gamma)=\widehat c\,|\gamma|\,u_L.
\]
La inversión conserva el número de pasos:
\[
 L_{\rm int}(\gamma^\vee)=L_{\rm int}(\gamma),
 \qquad
 L_{\leftrightarrow}=2L_{\rm int}(\gamma).
\]
Estas identidades son cinemáticas. Convertirlas en propagación
electromagnética exige que un mismo mapa reproduzca la ley de Gauss, la
dinámica de calibre, la dispersión nula y la reducción a dos grados
transversales.

\subsection{Entropía de entrelazamiento entre las \(2\) lecturas}
\label{sec:entrelazamiento-ida-retorno}

La inversión de una ruta no basta para afirmar entrelazamiento cuántico. Una
formulación mínima exige primero una completación de Hilbert. Sean
\(\mathcal H_{\rm adv}\) y \(\mathcal H_{\rm ret}\) dos espacios complejos
con familias ortonormales
\(\{|i\rangle_{\rm adv}\}_{i\in I}\) y
\(\{|i\rangle_{\rm ret}\}_{i\in I}\), donde \(I\) es finito. Para una
probabilidad \(p=(p_i)_{i\in I}\), definimos
\[
 |\Psi_p\rangle
 :=\sum_{i\in I}\sqrt{p_i}\,
 |i\rangle_{\rm adv}\otimes|i\rangle_{\rm ret}.
 \tag{E1}\label{eq:estado-entrelazado-ida-retorno}
\]

\begin{proposition}[Entropía bicapa de ida y retorno]
\label{prop:entropia-entrelazamiento-ida-retorno}
Las matrices reducidas del estado
\(\rho_p=|\Psi_p\rangle\langle\Psi_p|\) son
\[
 \rho_{\rm adv}
 =\sum_{i\in I}p_i|i\rangle_{\rm adv}\langle i|,
 \qquad
 \rho_{\rm ret}
 =\sum_{i\in I}p_i|i\rangle_{\rm ret}\langle i|,
 \tag{E2}\label{eq:reducidas-ida-retorno}
\]
y poseen la misma entropía de von Neumann:
\[
 \boxed{
 S_{\rm ent}(p)
 =-\operatorname{Tr}(\rho_{\rm adv}\log\rho_{\rm adv})
 =-\sum_{i\in I}p_i\log p_i.}
 \tag{E3}\label{eq:entropia-ida-retorno}
\]
Aquí se adopta la convención \(0\log0:=0\). El estado es separable
exactamente cuando una sola probabilidad vale uno.
\end{proposition}

\begin{proof}
Al efectuar la traza parcial, la ortogonalidad anula los términos con
\(i\ne j\). Los valores propios no nulos de cada matriz reducida son, por
tanto, los \(p_i\), lo que da \eqref{eq:entropia-ida-retorno}. El rango de
Schmidt es uno exactamente en el caso determinista.
\end{proof}

La proposición formaliza, una vez elegida la completación, una entropía entre
las lecturas avanzada y retardada. No deduce las probabilidades \(p_i\), no
identifica esas lecturas con los propagadores avanzado y retardado de una
teoría de campos y no convierte una correlación clásica de rutas en
entrelazamiento. La construcción de un isomorfismo que envíe la memoria HMT
a estados físicos bipartitos, preserve la evolución y fije \(p\) pertenece a
la \textsc{interfaz física}. La entropía nula de una distribución
determinista constituye el control negativo inmediato.

\begin{statusbox}
La proposición está
\textsc{demostrada con estructura de partida explícita}: presupone la
completación de Hilbert, las bases ortonormales y la distribución \(p\). La construcción
del entrelazador que realice físicamente ambas capas es una
\textsc{interfaz física}; no forma parte de la proposición.
\end{statusbox}

\begin{definition}[Interfaz fotónica de HMT--MD]
\label{def:interfaz-fotonica-hmt-md}
Una realización fotónica de HMT--MD es una sección abierta sostenida sobre
una dirección nula del álgebra partida, transportada entre fronteras
conjugadas y protegida por una simetría que conserva exactamente la norma
nula.
\end{definition}

Los idempotentes \(P_\pm\) no son, por sí solos, fotones físicos ni las dos
polarizaciones de Maxwell. La realización completa requiere un
entrelazador
\[
 \mathfrak R_\gamma:
 \partial_{\rm null}\mathcal C_{\rm sp}^{+}
 \dashrightarrow\mathcal H_{\rm Maxwell}^{\rm phys}
\]
que preserve el producto interno, la evolución y la acción de calibre, y que
seleccione exactamente dos polarizaciones transversales en \(3+1\)
dimensiones.

\begin{criterion}[Sector fotónico]
\label{crit:sector-fotonico}
La identificación física queda refutada si la evolución abandona
\(\partial_{\rm null}\mathcal C_{\rm sp}^{+}\), genera masa sin el mecanismo
de ruptura declarado, no reproduce la dispersión nula o no deja exactamente
dos polarizaciones físicas tras la reducción de calibre.
\end{criterion}

\section{Clausura material y centro electrónico}
\label{sec:clausura-material-electron}

El acoplamiento simultáneo de dos direcciones nulas conjugadas produce, para
\(u,v>0\),
\[
 N_{\rm sp}(uP_++vP_-)=uv>0.
\]
Ninguna componente aislada posee norma distinta de cero; la norma interior
surge de su acoplamiento. Esta afirmación algebraica no suma masas
preexistentes y no identifica todavía las direcciones nulas con fotones
\emph{on shell}.

   La clasificación estructural de HMT contiene un único punto fijo de la inversión celular,
\[
 (5,5).
\]
El origen relativo de acción del electrón es \(v_e=0\). En la realización
partida normalizada se fija
\[
 Z_e=P_++P_-=I,
 \qquad
 N_{\rm sp}(Z_e)=1.
\]
El punto \((5,5)\), el origen \(v_e=0\), la firma cruda
\((24,0,12,0,-12)\) y la palabra transversal \(555555\) son cuatro objetos
distintos.

\begin{definition}[Candidato de clausura electrónica]
\label{def:candidato-clausura-electronica}
La clausura resonante de modos nulos electromagnéticos conjugados se
formaliza como el programa que busca una sección material central normalizada
obtenida por el acoplamiento de dos modos nulos conjugados, estabilizada por
una clausura de ruta y por su holonomía.
\end{definition}

La existencia de la extensión, del retorno y del centro del atlas no prueba
la compactitud mínima ni la estabilidad del cierre. El teorema dinámico
pendiente debe construir
\[
 \gamma_e=\xi_1\circ\cdots\circ\xi_N,
 \qquad
 \partial\gamma_e=0,
 \qquad
 \varepsilon(\gamma_e)=-1,
\]
donde cada \(\xi_j\) sea localmente nulo, el borde total se anule y
\(\varepsilon\) sea la holonomía de signo. La masa se asociaría al cociclo
global de clausura, no a una suma de masas nulas locales. Este programa no
reabre el electrón central \((5,5)\), el atlas ni la cadena
extensión--ruta--registro--carácter ya construidos.

\begin{proposition}[Criterio relativista de interior causal]
\label{prop:cierre-nulos-interior-causal}
Sean \(p_1,\ldots,p_N\) momentos nulos futuros en un espacio de Minkowski con
firma temporal positiva, \(p_i^2=0\), y sea
\[
 P=\sum_{i=1}^{N}p_i.
\]
Si todos los \(p_i\) son colineales, entonces \(P^2=0\). Si al menos dos son
futuros y no colineales, entonces \(P^2>0\), y puede definirse
\[
 m^2c^2=P^2.
\]
\end{proposition}

\begin{proof}
Para vectores nulos futuros se tiene
\(p_i\cdot p_j\geq0\), con igualdad exactamente cuando son colineales.
Como
\[
 P^2=\sum_i p_i^2+2\sum_{i<j}p_i\cdot p_j,
\]
la primera suma es nula. Todos los productos cruzados se anulan en el caso
colineal; si existe un par no colineal, al menos uno es estrictamente
positivo.
\end{proof}

Por tanto, una ruta cerrada no es masiva por mera topología, ni una ruta
abierta es nula por el solo hecho de poseer borde. La realización exige que
el cierre geométrico y la suma de momentos entren conjuntamente en el
interior causal.

\subsection{Möbius, conjugación y ocupación posterior}

Bajo las hipótesis de levantamiento y holonomía del
\cref{thm:dos-cuatro-pi-estadistica}, una vuelta invierte la orientación y
dos la restauran. La banda de Möbius organiza partícula y antipartícula como
orientaciones conjugadas de un mismo tipo de cierre. Su identificación física
requiere una representación que invierta la carga, preserve la norma y
entrelace la evolución. La geometría de Möbius no implica por sí sola el
álgebra de ocupación fermiónica; sus relaciones operatorias se
introducen como hipótesis operatorias en el capítulo siguiente.

\begin{criterion}[Cierre electrónico]
\label{crit:cierre-electronico}
La realización queda refutada si el centro \((5,5)\) no se eleva a una
extensión cerrada, si la holonomía no selecciona el sector impar, si el
Hamiltoniano carece de una solución ligada estable y única salvo simetría, o
si carga, dispersión y factor de forma no concuerdan con la realización
electrónica.
\end{criterion}
